\documentclass[11pt,letterpaper]{article}

\usepackage[spanish,mexico,english]{babel}
\usepackage[utf8]{inputenc}
\usepackage[T1]{fontenc}
\usepackage[top=2.4cm,bottom=2.4cm,left=2.5cm,right=2.5cm]{geometry}
\usepackage{amsmath,amssymb}
\usepackage{graphicx}
\graphicspath{{../presentation/pictures/}}
\usepackage{booktabs}
\usepackage{float}
\usepackage{fancyhdr}
\usepackage{xcolor}
\usepackage{parskip}
\usepackage{hyperref}
\hypersetup{colorlinks=true,linkcolor=black,citecolor=black,urlcolor=blue}

\definecolor{uanlBlue}{RGB}{0,56,101}

\pagestyle{fancy}
\fancyhf{}
\renewcommand{\headrulewidth}{0.4pt}
\fancyhead[L]{\small\textcolor{uanlBlue}{MCIE -- Control No Lineal}}
\fancyhead[R]{\small\textcolor{uanlBlue}{Ene-Jun 2026}}
\fancyfoot[C]{\thepage}

\begin{document}

% ---------- PORTADA ----------
\thispagestyle{empty}
\begin{center}
  {\bfseries Universidad Autónoma de Nuevo León}\\
  Facultad de Ingeniería Mecánica y Eléctrica\\[0.6em]
  \textcolor{uanlBlue}{\rule{\linewidth}{1pt}}\\[0.6em]
  Maestría en Ciencias de la Ingeniería Eléctrica\\
  Unidad de Aprendizaje: \textbf{Control No Lineal}\\[1.2em]
  {\LARGE\bfseries Proyecto Integrador de Aprendizaje}\\[1em]
  \textcolor{uanlBlue}{\rule{0.6\linewidth}{0.4pt}}\\[1em]
  \begin{tabular}{rl}
    Estudiante: & Roberto Treviño Cervantes \\
    Matrícula:  & 1915003 \\
    Maestro:    & Dr.\ Alberto Cavazos \\
  \end{tabular}
\end{center}

\vfill
\newpage
\setcounter{page}{1}

\selectlanguage{english}

\section*{Case study}

\textit{Yield-Aware Higher-Order Trajectory Optimization for Step-and-Scan Wafer Stages.}

The commercial performance of a scanner is ultimately measured in wafers per hour. The time spent stepping and \textit{settling} --- the interval between the end of a step motion and the moment the stage has damped its residual vibrations to within the nanometre-level positioning specification --- represents a significant and controllable fraction of the total cycle time, which places the allowable settle time below 10\,ms. The key lever for reducing settling time is the generation of \textit{reference trajectories}: motion profiles that command the stage from its current position to the next scanning start-point with minimal excitation of mechanical resonances.

\section*{Model}

Al-Rawashdeh et al.\ propose modeling the lithography scanner as a combination of three open-loop serial chains: the reticle chain ($j=1$), the optics chain ($j=2$), and the wafer chain ($j=3$). A dual-stage architecture is employed: a long-stroke motor (LS) for large-range travel of at least 300\,mm with micrometer accuracy, and a short-stroke motor (SS) with Lorentz actuators for nanometric, contact-free positioning. The physical connections between stages are represented by lumped-parameter spring-damper systems according to the Kelvin-Voigt model.

The full 6-DOF model is restricted to in-plane motions, since rotations about $x$, $y$ and translation on $z$ are used mainly for focusing. The machine comprises 16 bodies; the global state vector has $6\times 16 = 96$ states.

\section*{Methodology}

Trajectories that have smooth higher-order derivatives result in less excitation of poorly damped modes in the structure. We implemented a simulation of the multi-body model using the MuJoCo physics engine for Python. The simulation serves as a virtual test bed where the proposed trajectories can be validated under the dynamic model.

The Kelvin-Voigt couplings are implemented as MuJoCo joints with explicit \texttt{stiffness} and \texttt{damping} parameters. For the actuated stages, MuJoCo position actuators are used with $K_P = 10^7$ at $dt = 10^{-4}$\,s. During the scanning phase, the wafer stage moves at $v_{w,\text{scan}} = 0.8$\,m/s while the reticle moves in the opposite direction at $v_{r,\text{scan}} = 3.2$\,m/s, maintaining the demagnification ratio $\alpha = 0.25$.

\section*{CNN cost function and validation}

The contribution of this work is the forward direction of the defect-prediction problem: given trajectory parameters $(v_{\max}, a_{\max}, j_{\max}, s_{\max})$, predict the expected wafer defect pattern by means of a CNN trained on the WM-811K dataset of real wafer maps. The cost function penalizes the cumulative probability mass on the non-plausible defect classes:
\begin{equation}
    J(\boldsymbol{r}) = \sum_{c\,\notin\,\mathcal{C}_{\text{scan}}} \hat{p}_c(\boldsymbol{r}),
    \qquad
    \mathcal{C}_{\text{scan}} = \{\textit{none}, \textit{scratch}, \textit{edge-local}, \textit{random}\}.
\end{equation}

The methodology is valid when synthetic wafer maps activate only the classes in $\mathcal{C}_{\text{scan}}$ and $J(\boldsymbol{\pi})$ shifts reproducibly under independent sweeps of $(v, a, j, s)_{\max}$. The CNN must never activate \textit{donut} or \textit{near-full}, which correspond to etching/cleaning patterns rather than trajectory-induced defects.

\selectlanguage{spanish}
\section*{Resultados}

\paragraph{Control PID + \textit{feedforward}.}
\begin{equation*}
    u(t) = \underbrace{K_p\, e(t) + K_i\!\int_0^{t}\! e(\tau)\,d\tau + K_d\,\frac{d\,e(t)}{dt}}_{\text{PID}} + \underbrace{m\,\ddot r(t)}_{\text{FF}}
\end{equation*}
{\footnotesize $e(t) = r(t) - y_{\text{larga}}(t)$\;\; (realimentación sobre la etapa larga).}

Ancho de banda $f_{bw}=50$ Hz, $\omega_c = 2\pi f_{bw}$, $\alpha\approx 3$.
\begin{center}
\footnotesize
\begin{tabular}{r@{$\;\approx\;$}l@{\qquad}r@{$\;=\;$}l}
$K_p$ & $8.66\times 10^{6}$ N/m            & $K_p$ & $m\,(2\pi f_{bw})^2$ \\
$K_i$ & $2.72\times 10^{8}$ N/(m$\cdot$s)  & $K_d$ & $m\,\omega_c$ \\
$K_d$ & $2.75\times 10^{4}$ N$\cdot$s/m    & $K_i$ & $K_p\,\omega_c/\alpha^2$ \\
$m$   & $87$ kg                            & \multicolumn{2}{l}{} \\
\end{tabular}
\end{center}

\begin{figure}[H]
    \centering
    \begin{minipage}{0.58\linewidth}\centering
        \includegraphics[width=\linewidth]{sim_combined_y.png}\\[2pt]
        {\footnotesize Error $e_{syn}$, MA y MSD (dir.\ Y).}
    \end{minipage}\hfill
    \begin{minipage}{0.38\linewidth}\centering
        \includegraphics[width=\linewidth]{scanning_3d.png}\\[2pt]
        {\footnotesize Entorno de simulación MuJoCo.}
    \end{minipage}
\end{figure}

\begin{figure}[H]
    \centering
    \includegraphics[width=0.55\linewidth]{cnn_confmat.png}\\[2pt]
    {\footnotesize Matriz de confusión --- 3 conv + 2 FC, 96.88\,\% \textit{accuracy}.}
\end{figure}

\selectlanguage{english}
\section*{Conclusion}

The forward direction explored here --- trajectory parameters $\to$ predicted defect pattern --- complements the established industrial inverse problem of inferring process causes from observed wafer defect maps. The CNN-based cost function enables trajectory optimization to proceed entirely in simulation, without access to physical lithographic equipment, and provides a reproducible link between the kinematic parameters of the reference trajectory and the predicted yield signature on the wafer.

\end{document}
